\section{Conclusion}
\label{sec:conclusion}

We asked whether an LLM, given only a formal grammar and the feedback a
black-box harness can actually produce -- acceptance rate, structural
diversity, and rejection signatures, with no coverage instrumentation at
all -- can iteratively refine a grammar-based fuzzing generator. Across five
seeded runs against a pinned cJSON build, it reliably can: mean acceptance
rate rose from 57.6\% to 96.8\% and structural diversity by 123\% relative
to a static baseline, a result with complete rank separation between arms
at the exact test's ceiling for this sample size. Reusing the same
pipeline, unmodified, against a second, independently adapted target
(parson) confirmed the mechanics generalize across targets, while five real
refinement iterations and 3{,}000 sanitizer-instrumented executions found no
crash -- a result we reported as an explainable negative finding, supported
by manual review of the highest-risk code paths, rather than reframed as a
success.

The broader point this project argues for is methodological rather than a
single number: a coverage-free proxy signal, cheap enough to compute from a
black-box harness's own accept/reject decision, is sufficient to steer
LLM-guided refinement of a grammar-based generator, which matters
specifically for the targets coverage-guided fuzzing cannot reach -- pinned
vendored dependencies, closed or license-restricted binaries, and embedded
or cross-compiled parsers. We also found that reconciling a formal grammar
against two real, independently implemented parsers is itself informative:
the same grammar under-specifies a directly opposite validity decision
(duplicate object keys accepted by one implementation, rejected by the
other) on two targets, which is exactly the kind of target-specific
knowledge a purely formal specification cannot supply and that fuzzing
against the real implementation surfaces for free.

We release the full pipeline, both harnesses, every campaign artifact, and
the reproducibility tooling described in Section~\ref{sec:methodology}, so
that the comparison in Section~\ref{sec:evaluation} and the null result on
parson can both be checked, extended, or falsified directly rather than
taken on faith.
